\documentclass[12pt,a4paper]{article}

% ---- Encoding & Font ----
\usepackage{iftex}
\ifPDFTeX
  \usepackage[utf8]{inputenc}
  \usepackage[T1]{fontenc}
\else
  \usepackage{fontspec}  % XeTeX/LuaTeX (tectonic): Unicode nativo (·, —, ¿, ª, §)
\fi
\usepackage[spanish]{babel}
\usepackage{csquotes}

% ---- Page layout ----
\usepackage[top=2.5cm, bottom=2.5cm, left=2.5cm, right=2.5cm]{geometry}
\usepackage{setspace}
\onehalfspacing

% ---- Math ----
\usepackage{amsmath, amssymb, amsthm}
\usepackage{mathtools}
\usepackage{physics}
\usepackage{esint}  % \oiint

% ---- Graphics ----
\usepackage{graphicx}
\usepackage{tikz}
\usepackage{float}
\usepackage{caption}

% ---- Tables ----
\usepackage{array}
\usepackage{booktabs}
\usepackage{tabularx}
\usepackage{colortbl}

% ---- Colours ----
\usepackage{xcolor}
\definecolor{notebg}{HTML}{FFF3CD}
\definecolor{noteborder}{HTML}{FFC107}
\definecolor{boxred}{HTML}{F8D7DA}
\definecolor{boxredborder}{HTML}{F5C6CB}

% ---- Boxes ----
\usepackage{mdframed}
\usepackage{tcolorbox}
\tcbuselibrary{most}

\newtcolorbox{keybox}[1][]{
  colback=blue!5,
  colframe=blue!70!black,
  arc=4pt,
  boxrule=1pt,
  fonttitle=\bfseries,
  title={#1}
}

\newtcolorbox{warnbox}[1][]{
  colback=boxred,
  colframe=boxredborder,
  coltitle=black!75,
  arc=4pt,
  boxrule=1pt,
  fonttitle=\bfseries,
  title={#1}
}

\newtcolorbox{notebox}[1][]{
  colback=notebg,
  colframe=noteborder,
  coltitle=black!75,
  arc=4pt,
  boxrule=1pt,
  fonttitle=\bfseries,
  title={#1}
}

% ---- Hyperlinks & TOC ----
\usepackage[hidelinks]{hyperref}
\usepackage{bookmark}

% ---- Enumerate ----
\usepackage{enumitem}

% ---- Miscellaneous ----
\usepackage{icomma}
\usepackage{siunitx}
\usepackage{textcomp}
\usepackage[bottom]{footmisc}

% ---- Diagramas (gráficas y esquemas) ----
\usepackage{pgfplots}
\pgfplotsset{compat=1.18}
\usetikzlibrary{babel}  % babel español activa `>`; esto evita romper las flechas `->`
% courses/ElecMag2024/Clases/assets/tikzstyles.tex
% ---------------------------------------------------------------------------
% Estilos compartidos para los diagramas del curso Electromagnetismo (2024).
% Fuente única del "look" visual (colores, grosores, escalas) para que todas
% las figuras (TikZ / pgfplots / CircuiTikz) sean consistentes entre clases.
%
% Es una COPIA de courses/Fisica3-2015/Clases/assets/tikzstyles.tex: mismo
% dominio, misma paleta. Copia y no symlink para que cada curso sea
% autocontenido y la paleta de Física III pueda evolucionar aparte.
%
% Uso: la clase debe cargar antes `tikz` y `pgfplots` (y `circuitikz` si la
%      clase tiene circuitos), y luego % courses/ElecMag2024/Clases/assets/tikzstyles.tex
% ---------------------------------------------------------------------------
% Estilos compartidos para los diagramas del curso Electromagnetismo (2024).
% Fuente única del "look" visual (colores, grosores, escalas) para que todas
% las figuras (TikZ / pgfplots / CircuiTikz) sean consistentes entre clases.
%
% Es una COPIA de courses/Fisica3-2015/Clases/assets/tikzstyles.tex: mismo
% dominio, misma paleta. Copia y no symlink para que cada curso sea
% autocontenido y la paleta de Física III pueda evolucionar aparte.
%
% Uso: la clase debe cargar antes `tikz` y `pgfplots` (y `circuitikz` si la
%      clase tiene circuitos), y luego % courses/ElecMag2024/Clases/assets/tikzstyles.tex
% ---------------------------------------------------------------------------
% Estilos compartidos para los diagramas del curso Electromagnetismo (2024).
% Fuente única del "look" visual (colores, grosores, escalas) para que todas
% las figuras (TikZ / pgfplots / CircuiTikz) sean consistentes entre clases.
%
% Es una COPIA de courses/Fisica3-2015/Clases/assets/tikzstyles.tex: mismo
% dominio, misma paleta. Copia y no symlink para que cada curso sea
% autocontenido y la paleta de Física III pueda evolucionar aparte.
%
% Uso: la clase debe cargar antes `tikz` y `pgfplots` (y `circuitikz` si la
%      clase tiene circuitos), y luego \input{../assets/tikzstyles.tex}
% (compilar desde el directorio de la clase para que la ruta relativa resuelva).
% ---------------------------------------------------------------------------

% ---- Paleta (alineada con las cajas del preámbulo: keybox/notebox/warnbox) ----
\definecolor{figblue}{HTML}{1F4E79}   % azul principal (líneas, keybox)
\definecolor{figred}{HTML}{C0392B}    % rojo de énfasis (warnbox)
\definecolor{figamber}{HTML}{B8860B}  % ámbar (notebox)
\definecolor{figgray}{HTML}{555555}   % auxiliares, ejes, guías

% ---- CircuiTikz: escala y grosor de los componentes ----
% Guarda \ifdefined: la electrostática (Clases 1-14) no carga `circuitikz`, y
% sin la guarda el \input revienta con `Undefined control sequence \ctikzset`.
% Cuando el curso llegue a circuitos, el bloque se activa solo.
\ifdefined\ctikzset
  \ctikzset{
    bipoles/thickness=1,
    resistors/scale=0.9,
    capacitors/scale=0.9,
  }
\fi

% ---- TikZ: estilos reutilizables ----
% Nota: las puntas se declaran como `-{Latex[...]}` y no como `->`. La forma
% con llaves NO contiene un `>` literal, así que es inmune al `>` activo que
% introduce babel español (ver CLAUDE.md §7.3.3).
\usetikzlibrary{arrows.meta,calc,angles,quotes}

\tikzset{
  figline/.style={figblue, line width=0.9pt},
  figaux/.style={figgray, dashed, line width=0.6pt},
  % vectores
  figvec/.style={figblue, line width=0.9pt, -{Latex[length=2.2mm,width=1.6mm]}},
  figvecres/.style={figred, line width=1.3pt, -{Latex[length=2.6mm,width=1.9mm]}},
  figvecaux/.style={figgray, line width=0.7pt, -{Latex[length=1.8mm,width=1.3mm]}},
  figaxis/.style={figgray, line width=0.6pt, -{Latex[length=2mm,width=1.4mm]}},
  % cargas puntuales
  figcarga/.style={circle, inner sep=0pt, minimum size=5.4pt, draw=none},
  figpos/.style={figcarga, fill=figred},
  figneg/.style={figcarga, fill=figblue},
  figneutra/.style={figcarga, fill=figgray},
  % etiquetas
  figlbl/.style={font=\small},
  figlblaux/.style={font=\footnotesize, figgray},
}

% ---- pgfplots: estilo base de las gráficas ----
\pgfplotsset{
  emfig/.style={
    width=11cm, height=6.5cm,
    axis lines=left,
    xlabel style={font=\small},
    ylabel style={font=\small},
    tick label style={font=\footnotesize},
    every axis plot/.append style={line width=1pt},
    grid=major,
    grid style={figgray!20},
    legend cell align=left,
    legend style={font=\footnotesize, draw=figgray!40},
  },
}

% (compilar desde el directorio de la clase para que la ruta relativa resuelva).
% ---------------------------------------------------------------------------

% ---- Paleta (alineada con las cajas del preámbulo: keybox/notebox/warnbox) ----
\definecolor{figblue}{HTML}{1F4E79}   % azul principal (líneas, keybox)
\definecolor{figred}{HTML}{C0392B}    % rojo de énfasis (warnbox)
\definecolor{figamber}{HTML}{B8860B}  % ámbar (notebox)
\definecolor{figgray}{HTML}{555555}   % auxiliares, ejes, guías

% ---- CircuiTikz: escala y grosor de los componentes ----
% Guarda \ifdefined: la electrostática (Clases 1-14) no carga `circuitikz`, y
% sin la guarda el \input revienta con `Undefined control sequence \ctikzset`.
% Cuando el curso llegue a circuitos, el bloque se activa solo.
\ifdefined\ctikzset
  \ctikzset{
    bipoles/thickness=1,
    resistors/scale=0.9,
    capacitors/scale=0.9,
  }
\fi

% ---- TikZ: estilos reutilizables ----
% Nota: las puntas se declaran como `-{Latex[...]}` y no como `->`. La forma
% con llaves NO contiene un `>` literal, así que es inmune al `>` activo que
% introduce babel español (ver CLAUDE.md §7.3.3).
\usetikzlibrary{arrows.meta,calc,angles,quotes}

\tikzset{
  figline/.style={figblue, line width=0.9pt},
  figaux/.style={figgray, dashed, line width=0.6pt},
  % vectores
  figvec/.style={figblue, line width=0.9pt, -{Latex[length=2.2mm,width=1.6mm]}},
  figvecres/.style={figred, line width=1.3pt, -{Latex[length=2.6mm,width=1.9mm]}},
  figvecaux/.style={figgray, line width=0.7pt, -{Latex[length=1.8mm,width=1.3mm]}},
  figaxis/.style={figgray, line width=0.6pt, -{Latex[length=2mm,width=1.4mm]}},
  % cargas puntuales
  figcarga/.style={circle, inner sep=0pt, minimum size=5.4pt, draw=none},
  figpos/.style={figcarga, fill=figred},
  figneg/.style={figcarga, fill=figblue},
  figneutra/.style={figcarga, fill=figgray},
  % etiquetas
  figlbl/.style={font=\small},
  figlblaux/.style={font=\footnotesize, figgray},
}

% ---- pgfplots: estilo base de las gráficas ----
\pgfplotsset{
  emfig/.style={
    width=11cm, height=6.5cm,
    axis lines=left,
    xlabel style={font=\small},
    ylabel style={font=\small},
    tick label style={font=\footnotesize},
    every axis plot/.append style={line width=1pt},
    grid=major,
    grid style={figgray!20},
    legend cell align=left,
    legend style={font=\footnotesize, draw=figgray!40},
  },
}

% (compilar desde el directorio de la clase para que la ruta relativa resuelva).
% ---------------------------------------------------------------------------

% ---- Paleta (alineada con las cajas del preámbulo: keybox/notebox/warnbox) ----
\definecolor{figblue}{HTML}{1F4E79}   % azul principal (líneas, keybox)
\definecolor{figred}{HTML}{C0392B}    % rojo de énfasis (warnbox)
\definecolor{figamber}{HTML}{B8860B}  % ámbar (notebox)
\definecolor{figgray}{HTML}{555555}   % auxiliares, ejes, guías

% ---- CircuiTikz: escala y grosor de los componentes ----
% Guarda \ifdefined: la electrostática (Clases 1-14) no carga `circuitikz`, y
% sin la guarda el \input revienta con `Undefined control sequence \ctikzset`.
% Cuando el curso llegue a circuitos, el bloque se activa solo.
\ifdefined\ctikzset
  \ctikzset{
    bipoles/thickness=1,
    resistors/scale=0.9,
    capacitors/scale=0.9,
  }
\fi

% ---- TikZ: estilos reutilizables ----
% Nota: las puntas se declaran como `-{Latex[...]}` y no como `->`. La forma
% con llaves NO contiene un `>` literal, así que es inmune al `>` activo que
% introduce babel español (ver CLAUDE.md §7.3.3).
\usetikzlibrary{arrows.meta,calc,angles,quotes}

\tikzset{
  figline/.style={figblue, line width=0.9pt},
  figaux/.style={figgray, dashed, line width=0.6pt},
  % vectores
  figvec/.style={figblue, line width=0.9pt, -{Latex[length=2.2mm,width=1.6mm]}},
  figvecres/.style={figred, line width=1.3pt, -{Latex[length=2.6mm,width=1.9mm]}},
  figvecaux/.style={figgray, line width=0.7pt, -{Latex[length=1.8mm,width=1.3mm]}},
  figaxis/.style={figgray, line width=0.6pt, -{Latex[length=2mm,width=1.4mm]}},
  % cargas puntuales
  figcarga/.style={circle, inner sep=0pt, minimum size=5.4pt, draw=none},
  figpos/.style={figcarga, fill=figred},
  figneg/.style={figcarga, fill=figblue},
  figneutra/.style={figcarga, fill=figgray},
  % etiquetas
  figlbl/.style={font=\small},
  figlblaux/.style={font=\footnotesize, figgray},
}

% ---- pgfplots: estilo base de las gráficas ----
\pgfplotsset{
  emfig/.style={
    width=11cm, height=6.5cm,
    axis lines=left,
    xlabel style={font=\small},
    ylabel style={font=\small},
    tick label style={font=\footnotesize},
    every axis plot/.append style={line width=1pt},
    grid=major,
    grid style={figgray!20},
    legend cell align=left,
    legend style={font=\footnotesize, draw=figgray!40},
  },
}


% ---- Title ----
\title{\textbf{Clase 9: Polarización, cargas ligadas y ley de Gauss en dieléctricos} \\[0.3em]
        \large{El vector polarización, las cargas de polarización, el campo macroscópico y el vector desplazamiento}}
\author{Transcripto y expandido --- Curso Electromagnetismo (OpenFING)}
\date{}

\begin{document}

\maketitle
\thispagestyle{empty}
\newpage

\tableofcontents
\newpage

\section{Por qué un aislante sí importa}

Hasta ahora todo el curso transcurrió en el vacío, con la única excepción de los
conductores. Empieza el capítulo 3: qué pasa cuando hay \textbf{materiales
dieléctricos}.

Un dieléctrico, por definición, \textbf{no contiene cargas libres}: no tiene
cargas que puedan moverse distancias macroscópicas de un lado a otro. La reacción
inmediata sería concluir que entonces es electrostáticamente irrelevante. Es
falsa, y el motivo es que las cargas sí pueden moverse \textbf{distancias
microscópicas}.

Un plástico está formado por un número enorme de moléculas y, frecuentemente,
cada una de ellas está polarizada: el baricentro de las cargas positivas no
coincide con el de las negativas, de modo que la molécula tiene un \textbf{momento
dipolar} propio y una dirección preferencial.

Al aplicar un campo eléctrico externo, cada dipolo tiende a orientarse: la carga
negativa quiere acercarse a las cargas positivas que generaron el campo, y la
positiva a las negativas. El resultado no es un alineamiento perfecto ---las
moléculas siguen con su agitación térmica--- pero \textbf{en promedio} el efecto
se acumula, y esa acumulación sobre un número enorme de moléculas produce un
efecto macroscópico: la \textbf{polarización}.

\begin{warnbox}[La figura de las moléculas alineadas es mentira]
El docente insiste en que el dibujo de todas las moléculas alineadas es una
\textbf{caricatura}. En la realidad las moléculas siguen casi tan desordenadas
como antes; están sólo \emph{levemente} orientadas. Lo que ocurre es que sin
campo el promedio da cero y con campo no.
\end{warnbox}

\begin{figure}[H]
  \centering
  % Clase 9 — la caricatura de la orientacion molecular, y su desmentida.
% El docente dibuja las moleculas alineadas y ACTO SEGUIDO aclara que es
% mentira ("ojo, ojo, imagen, caricatura"). La figura tiene que mostrar las dos
% cosas o repite el error que el se preocupa por corregir: de ahi el tercer
% panel, que es el que dice la verdad.
% Los dipolos se dibujan como un par de puntos unidos; el angulo de cada uno se
% fija a mano y no al azar, para que el panel "realista" se vea desordenado pero
% con un sesgo leve y reproducible.
\begin{tikzpicture}[scale=1]

  % ---------- panel 1: sin campo ----------
  \begin{scope}[shift={(0,0)}]
    \draw[figgray!45, line width=0.6pt] (-1.5,-1.5) rectangle (1.5,1.5);
    \foreach \x/\y/\a in {-1.0/1.0/25, -0.1/1.05/200, 0.9/0.95/110,
                          -1.05/0.1/300, -0.05/0.0/70, 1.0/0.15/160,
                          -0.95/-1.0/135, 0.0/-1.05/340, 1.05/-0.95/55}{
      \begin{scope}[shift={(\x,\y)}, rotate=\a]
        \draw[figline, line width=0.7pt] (-0.22,0) -- (0.22,0);
        \node[figneg, minimum size=4pt] at (-0.22,0) {};
        \node[figpos, minimum size=4pt] at (0.22,0) {};
      \end{scope}}
    \node[figlbl] at (0,-1.95) {sin campo};
    \node[figlblaux, align=center] at (0,-2.6) {orientaciones\\ aleatorias};
  \end{scope}

  % ---------- panel 2: la caricatura ----------
  \begin{scope}[shift={(4.4,0)}]
    \draw[figgray!45, line width=0.6pt] (-1.5,-1.5) rectangle (1.5,1.5);
    \foreach \x in {-1.0,0.0,1.0}{
      \foreach \y in {-1.0,0.0,1.0}{
        \begin{scope}[shift={(\x,\y)}, rotate=90]
          \draw[figline, line width=0.7pt] (-0.22,0) -- (0.22,0);
          \node[figneg, minimum size=4pt] at (-0.22,0) {};
          \node[figpos, minimum size=4pt] at (0.22,0) {};
        \end{scope}}}
    \node[figlbl] at (0,-1.95) {la caricatura};
    \node[figlblaux, figred, align=center] at (0,-2.6) {\emph{no} es lo que pasa};
  \end{scope}

  % ---------- panel 3: lo que realmente pasa ----------
  \begin{scope}[shift={(8.8,0)}]
    \draw[figgray!45, line width=0.6pt] (-1.5,-1.5) rectangle (1.5,1.5);
    % mismos angulos que el panel 1, sesgados ~15 grados hacia la vertical
    \foreach \x/\y/\a in {-1.0/1.0/40, -0.1/1.05/195, 0.9/0.95/105,
                          -1.05/0.1/295, -0.05/0.0/85, 1.0/0.15/150,
                          -0.95/-1.0/125, 0.0/-1.05/335, 1.05/-0.95/70}{
      \begin{scope}[shift={(\x,\y)}, rotate=\a]
        \draw[figline, line width=0.7pt] (-0.22,0) -- (0.22,0);
        \node[figneg, minimum size=4pt] at (-0.22,0) {};
        \node[figpos, minimum size=4pt] at (0.22,0) {};
      \end{scope}}
    \node[figlbl] at (0,-1.95) {lo que pasa};
    \node[figlblaux, align=center] at (0,-2.6) {casi igual de desordenadas,\\ apenas sesgadas};
  \end{scope}

  % ---------- campo aplicado, comun a los paneles 2 y 3 ----------
  \foreach \x in {3.2,5.6,7.6,10.0}{
    \draw[figvecaux] (\x,-1.5) -- (\x,1.5);}
  \node[figlbl, figgray] at (6.6,1.85) {$\vec E$ aplicado};

  \node[figlbl, align=center] at (4.4,-3.6)
    {El promedio del panel 1 da cero; el del panel 3, no. \textbf{Ese} promedio no nulo};
  \node[figlbl, align=center] at (4.4,-4.15)
    {es la polarizaci\'on --- no el alineamiento perfecto del panel 2.};

\end{tikzpicture}

  \caption{La caricatura y su desmentida. El efecto macroscópico no viene del
  alineamiento perfecto sino de un sesgo minúsculo sobre un número enorme de
  moléculas.}
\end{figure}

\subsection{Moléculas polares y apolares}

También existen moléculas \textbf{apolares}, sin momento dipolar propio. Todo lo
que sigue vale igual para ellas, pero por una razón más sofisticada: el propio
campo aplicado deforma la molécula y le genera un \textbf{dipolo inducido}. El
efecto es mucho más chico y el tratamiento más complicado.

\begin{notebox}[Alcance]
La clase se restringe a \textbf{moléculas polares}, tomando los dipolos como dato
y sin considerar la deformación de la molécula por el campo aplicado.
\end{notebox}

\subsection{De la molécula al dipolo puntual}

Una molécula real no es un dipolo: es neutra, pero tiene momentos cuadrupolares,
octopolares y todo lo demás. La simplificación viene de que las moléculas son
\textbf{chiquitas}: su tamaño típico es el nanómetro, y cualquier distancia
macroscópica ---aunque sea una micra--- ya es enormemente mayor.

Por el desarrollo multipolar, un objeto de \textbf{carga neta nula} visto a
distancias mucho mayores que su tamaño se comporta como un \textbf{dipolo}: los
demás momentos decaen más rápido. Entonces, a efectos macroscópicos, cada
molécula puede reemplazarse por su momento dipolar y toda la demás información se
descarta.

\section{El vector polarización}

Lo que le pasa a cada molécula individual no interesa: a escala molecular el
problema es intratable. Lo que interesa es el \textbf{efecto promediado} de
muchas moléculas en un volumen pequeño.

Se toma un volumen $\Delta V$ \textbf{grande respecto de la escala molecular y
pequeño respecto de la macroscópica} ---por ejemplo un cubo de una micra de lado
dentro de un objeto de \SI{20}{\centi\metre}---, que contiene una cantidad enorme
de moléculas. Su momento dipolar total es
\[
\Delta \vec p = \sum_{i \in \Delta V} \vec p_i,
\qquad
\vec p_i = \int_{\text{molécula}} \rho(\vec r_i)\, \vec r_i \, dV .
\]

Pero $\Delta\vec p$ no sirve como magnitud local: en un material homogéneo, al
duplicar el volumen se duplica el momento dipolar. Lo que sí tiene límite es la
densidad.

\begin{keybox}[Campo de polarización]
\[
\boxed{\ \vec P(\vec r) = \lim_{\Delta V \to 0} \frac{\Delta \vec p}{\Delta V}\ }
\]
el momento dipolar por unidad de volumen.
\end{keybox}

\begin{warnbox}[El límite es un abuso de notación]
$\Delta V \to 0$ significa pequeño \textbf{macroscópicamente}, nunca
microscópicamente. A escala molecular la polarización deja de tener sentido,
porque se pone a variar de molécula a molécula.
\end{warnbox}

Si el material es homogéneo, $\vec P$ no depende de la posición; si es
inhomogéneo, varía suavemente a escala macroscópica.

\section{Potencial de un dieléctrico polarizado}

El problema que se plantea es el inverso del que uno esperaría: \textbf{la
polarización se da como dato}. No se pregunta todavía cómo se polarizó el
material, sino qué campo eléctrico produce una polarización conocida.

Se considera un material que ocupa un volumen $V_0$, con polarización
$\vec P(\vec r\,')$, y se observa el potencial en un punto $\vec r$ \textbf{fuera
del material}. Cada elemento de volumen $dV'$ tiene momento dipolar
$\vec P(\vec r\,')\,dV'$ y aporta un potencial dipolar. Sumando:

\begin{keybox}[Potencial generado por una polarización conocida]
\[
\boxed{\ \phi(\vec r) = \frac{1}{4\pi\varepsilon_0}
\int_{V_0} \frac{\vec P(\vec r\,') \cdot (\vec r - \vec r\,')}
{|\vec r - \vec r\,'|^3}\, dV' \ }
\]
\end{keybox}

\begin{figure}[H]
  \centering
  % Clase 9 — la geometria del potencial de un dielectrico polarizado.
% Es el dibujo que el docente rehace al empezar la clase ("teniamos un volumen
% B0 ... el origen aca y un vector R'"). Lo que la figura tiene que dejar claro
% es CUAL de los tres vectores se integra: r' recorre el material, r esta fijo.
% El punto de observacion se pone bien afuera y con una nota: la deduccion
% empieza valiendo solo fuera del material, y recien la cavidad en forma de
% aguja la extiende hacia adentro.
\begin{tikzpicture}[scale=1]

  % cuerpo dielectrico, forma irregular a proposito (no es una esfera)
  \draw[figline, fill=figblue!8]
    plot[smooth cycle, tension=0.8]
    coordinates {(-2.5,0.2) (-1.7,1.5) (-0.2,1.8) (1.0,1.1) (0.9,-0.4)
                 (-0.1,-1.3) (-1.6,-1.2)};
  \node[figlbl, figblue] at (-1.9,1.15) {$V_0$};
  \node[figlblaux] at (-0.55,-1.75) {superficie $S_0$};

  % origen
  \node[figneutra] at (-3.6,-2.3) {};
  \node[figlblaux] at (-3.95,-2.5) {$O$};

  % elemento de volumen y su dipolo
  \coordinate (fuente) at (-0.9,0.55);
  \fill[figblue!30] (fuente) circle (0.11);
  \node[figlblaux, align=center] at (-0.95,1.02) {$dV'$};
  \draw[figvec, figred, line width=0.8pt] ($(fuente)+(-0.28,-0.18)$) -- ($(fuente)+(0.28,0.18)$);
  \node[figlbl, figred] at (-0.15,0.38) {$\vec P(\vec r\,')$};

  % punto de observacion
  \node[figneutra] at (3.3,1.35) {};
  \node[figlbl] at (3.62,1.55) {$\vec r$};

  % los tres vectores
  \draw[figvecaux] (-3.6,-2.3) -- (fuente);
  \node[figlbl, figgray, fill=white, inner sep=1.2pt] at (-2.55,-0.95) {$\vec r\,'$};

  \draw[figvec] (-3.6,-2.3) -- (3.3,1.35);
  \node[figlbl, figblue, fill=white, inner sep=1.2pt] at (1.35,-0.95) {$\vec r$};

  \draw[figvecres] (fuente) -- (3.3,1.35);
  \node[figlbl, figred, fill=white, inner sep=1.2pt] at (1.5,1.35) {$\vec r-\vec r\,'$};

  \node[figlblaux, align=center] at (0.4,-2.9)
    {$\vec r$ \emph{fuera} de $V_0$: por ahora la deducci\'on vale s\'olo ah\'i.};

\end{tikzpicture}

  \caption{Se integra sobre $\vec r\,'$, que recorre el material; el punto de
  observación $\vec r$ queda fijo. La deducción vale, por ahora, sólo con
  $\vec r$ fuera de $V_0$.}
\end{figure}

\begin{notebox}[Por qué la aproximación dipolar es legítima para todas]
Estando $\vec r$ fuera del material, incluso la molécula más cercana está a una
distancia macroscópica, enormemente mayor que su propio tamaño.
\end{notebox}

Este es sólo el potencial \textbf{producido por la polarización}. Si además hay
cargas libres ---las que uno deposita o manipula--- sus potenciales se suman
aparte.

\subsection{El truco del gradiente primado}

Se define $\nabla'$ como el gradiente respecto de las coordenadas de
$\vec r\,'$. La observación clave es
\[
\nabla' \frac{1}{|\vec r - \vec r\,'|}
= \frac{\vec r - \vec r\,'}{|\vec r - \vec r\,'|^3},
\]
que es el mismo objeto que aparece en el campo de una carga puntual, salvo el
signo que introduce derivar respecto de la variable primada en lugar de la no
primada. Con eso el integrando se escribe como
$\vec P \cdot \nabla'\!\left(1/|\vec r - \vec r\,'|\right)$.

\begin{notebox}[Por qué el gradiente primado y no el otro]
Para aplicar el teorema de la divergencia hacen falta derivadas respecto de la
\textbf{misma} variable sobre la que se integra, no respecto de una variable que
queda afuera de la integral.
\end{notebox}

Eso todavía no es una divergencia, pero se parece. Usando la identidad
\[
\div(f\vec A) = \grad f \cdot \vec A + f \, \div \vec A
\]
leída al revés, el integrando se parte en dos:
\[
\phi(\vec r) = \frac{1}{4\pi\varepsilon_0} \int_{V_0}
\left[ \nabla' \cdot \frac{\vec P(\vec r\,')}{|\vec r - \vec r\,'|}
- \frac{\nabla' \cdot \vec P(\vec r\,')}{|\vec r - \vec r\,'|} \right] dV' .
\]

Al primer término ---y sólo a ese--- se le aplica el teorema de la divergencia,
que lo convierte en una integral sobre la superficie $S_0$ que bordea $V_0$:
\[
\phi(\vec r) = \frac{1}{4\pi\varepsilon_0} \oiint_{S_0}
\frac{\vec P(\vec r\,') \cdot \hat n'}{|\vec r - \vec r\,'|}\, dS'
+ \frac{1}{4\pi\varepsilon_0} \int_{V_0}
\frac{-\nabla' \cdot \vec P(\vec r\,')}{|\vec r - \vec r\,'|}\, dV' .
\]

\subsection{Las cargas de polarización}

Las dos integrales tienen exactamente la forma coulombiana del potencial de una
densidad superficial y de una densidad volumétrica de carga. Eso permite
\textbf{bautizar} lo que aparece arriba.

\begin{keybox}[Densidades de carga de polarización]
\[
\boxed{\ \rho_p(\vec r) = -\div \vec P
\qquad\text{y}\qquad
\sigma_p(\vec r) = \vec P \cdot \hat n \ }
\]
con $\hat n$ \textbf{saliente del material}. La orientación queda bien definida
porque la polarización está asociada a un material, que tiene un adentro y un
afuera.
\end{keybox}

Con estos nombres el potencial queda idéntico al de cualquier distribución de
carga:
\[
\phi(\vec r) = \frac{1}{4\pi\varepsilon_0} \oiint_{S_0}
\frac{\sigma_p(\vec r\,')}{|\vec r - \vec r\,'|}\, dS'
+ \frac{1}{4\pi\varepsilon_0} \int_{V_0}
\frac{\rho_p(\vec r\,')}{|\vec r - \vec r\,'|}\, dV' .
\]

\begin{notebox}[Nombre y apellido]
Estas cargas «tienen nombre carga y apellido polarización»: si uno se olvida del
apellido, la fórmula es la usual. La única particularidad es que $\rho_p$ y
$\sigma_p$ \textbf{no son arbitrarias}, sino que se expresan de manera precisa en
términos de $\vec P$.
\end{notebox}

El campo eléctrico se obtiene entonces sin ninguna complicación adicional, porque
\textbf{Coulomb sigue valiendo para todas las cargas}:
\[
\vec E(\vec r) = \frac{1}{4\pi\varepsilon_0} \int_{V_0}
\rho_p(\vec r\,') \frac{\vec r - \vec r\,'}{|\vec r - \vec r\,'|^3}\, dV'
+ \frac{1}{4\pi\varepsilon_0} \oiint_{S_0}
\sigma_p(\vec r\,') \frac{\vec r - \vec r\,'}{|\vec r - \vec r\,'|^3}\, dS' .
\]

Las cargas de polarización son producto de la reorientación local de las
moléculas, pero no por eso dejan de ser cargas.

\section{Por qué siempre aparece carga superficial}

Que aparezca una densidad \textbf{superficial} en un dieléctrico puede
sorprender: es el tipo de cosa que uno asocia a conductores. La razón es que en
el borde hay un \textbf{cambio abrupto}.

El ejemplo es un condensador con carga libre positiva en una placa y un
dieléctrico en el medio. Con campo aplicado uniforme, las moléculas se orientan
todas igual (caricatura mediante): en cualquier región interior hay tantas cargas
positivas como negativas y \textbf{no queda densidad volumétrica}, pero sobra un
excedente de carga negativa de un lado y positiva del otro.

\begin{figure}[H]
  \centering
  % Clase 9 — por que aparece carga superficial de polarizacion y no volumetrica.
% Es el ejemplo del condensador con dielectrico. El contenido de la figura es la
% CANCELACION: en el interior cada carga negativa de un dipolo queda pegada a la
% positiva del siguiente, y lo unico que sobra son las dos caras.
% Los dipolos se dibujan en columnas alineadas a proposito (es la caricatura del
% campo uniforme) y el sombreado marca la region donde todo se compensa.
\begin{tikzpicture}[scale=1]

  % placas del condensador
  \draw[figline, line width=1.6pt, figred] (-3.2,2.1) -- (3.2,2.1);
  \draw[figline, line width=1.6pt, figblue] (-3.2,-2.1) -- (3.2,-2.1);
  \foreach \x in {-2.6,-1.3,0,1.3,2.6}{
    \node[figlbl, figred] at (\x,2.42) {$+$};
    \node[figlbl, figblue] at (\x,-2.45) {$-$};}
  \node[figlblaux, figred] at (-4.0,2.1) {carga libre};
  \node[figlblaux, figblue] at (-4.0,-2.1) {carga libre};

  % dielectrico
  \fill[figamber!10] (-2.9,-1.55) rectangle (2.9,1.55);
  \draw[figamber, line width=0.8pt] (-2.9,-1.55) rectangle (2.9,1.55);

  % zona de cancelacion
  \fill[figgray!10] (-2.6,-1.05) rectangle (2.6,1.05);

  % dipolos alineados (caricatura del campo uniforme)
  \foreach \x in {-2.2,-1.1,0,1.1,2.2}{
    \foreach \y in {-1.0,0.0,1.0}{
      \begin{scope}[shift={(\x,\y)}]
        \draw[figline, line width=0.7pt] (0,-0.3) -- (0,0.3);
        \node[figneg, minimum size=4.4pt] at (0,-0.3) {};
        \node[figpos, minimum size=4.4pt] at (0,0.3) {};
      \end{scope}}}

  % lo que sobra en las caras
  \node[figlblaux, figblue, fill=white, inner sep=1.4pt] at (0,1.33) {$\sigma_p<0$};
  \node[figlblaux, figred, fill=white, inner sep=1.4pt] at (0,-1.33) {$\sigma_p>0$};
  % Nada de rotular rho_p=0 aca: el unico hueco libre es donde va la flecha de
  % E, y el caption ya lo dice. La figura muestra la cancelacion; el texto la
  % nombra.

  % campo aplicado
  \draw[figvecaux] (3.75,1.4) -- (3.75,-1.4);
  \node[figlbl, figgray] at (4.15,0) {$\vec E$};

\end{tikzpicture}

  \caption{Con campo aplicado uniforme, en el interior cada carga negativa queda
  pegada a la positiva de su vecina y todo se compensa: no hay densidad
  volumétrica y sólo sobran las dos caras. Con campo inhomogéneo aparecen las
  dos densidades.}
\end{figure}

\begin{table}[H]
  \centering
  \begin{tabular}{lcc}
    \toprule
    Campo aplicado & $\rho_p$ & $\sigma_p$ \\
    \midrule
    Uniforme      & $0$                   & $\neq 0$ \\
    Inhomogéneo   & en general $\neq 0$   & en general $\neq 0$ \\
    \bottomrule
  \end{tabular}
  \caption{El caso uniforme, con sólo carga superficial, es el que se veía en
  Física III. Ahora aparecen las dos, y la superficial no se puede olvidar.}
\end{table}

\section{La carga de polarización total es cero}

Tiene que serlo: la polarización es producto de dipolos y cada dipolo tiene carga
neta nula. Pero conviene verificarlo, porque si a uno le dieran un $\rho$ y un
$\sigma$ cualesquiera no habría motivo para que la suma se anulara. Que se anule
es consecuencia de que \textbf{no son cualesquiera}.
\[
Q_p = \int_{V_0} \rho_p \, dV' + \oiint_{S_0} \sigma_p \, dS'
= \int_{V_0} \left(-\div' \vec P\right) dV'
+ \oiint_{S_0} \vec P \cdot \hat n'\, dS' .
\]

Aplicando el teorema de la divergencia al primer término se obtiene
$-\oiint_{S_0} \vec P \cdot \hat n'\, dS'$, que cancela exactamente al segundo:

\begin{keybox}[Neutralidad de la polarización]
\[
\boxed{\ Q_p = 0\ }
\]
\end{keybox}

\begin{notebox}[Regla mnemotécnica]
Uno se puede olvidar cuál de las dos definiciones lleva el signo menos, pero
\textbf{una lleva más y la otra menos}, y es justamente por esta cancelación.
\end{notebox}

\section{El campo macroscópico}

Todo lo anterior se dedujo observando desde \textbf{fuera} del material. Adentro
hay que ser mucho más cuidadoso con \emph{qué} campo eléctrico se está nombrando.

El campo \textbf{microscópico} dentro del material varía salvajemente: cerca de
una carga positiva se dispara, cerca de una negativa se dispara para el otro
lado, y además depende fuertemente del tiempo porque las moléculas se agitan sin
parar. No es ese el campo de interés.

El campo de interés es el que uno \textbf{mediría con un aparato}. Un voltímetro
con sus agujas metidas en el material da una diferencia de potencial estable, que
no oscila: la punta de la aguja, aunque sea fina macroscópicamente, abarca una
cantidad enorme de moléculas y \textbf{promedia} sobre todas ellas.

Formalmente se define, con una función de peso $f$,
\[
\vec E_{\text{macro}}(\vec r) = \int_V f(\vec r\,')\,
\vec E_{\text{micro}}(\vec r + \vec r\,')\, dV',
\qquad \int_V f(\vec r\,')\, dV' = 1,
\]
donde $f$ está concentrada cerca de $\vec r$ y decae al alejarse, y $V$ es grande
frente a la escala molecular y pequeño frente a la macroscópica. El peso total es
1 para no sesgar el promedio.

\begin{notebox}[El detalle del promedio no importa]
Se puede tomar $f$ uniforme ($1/V$), pero tiene el defecto del salto abrupto en
el borde; conviene una función suave. Otra opción equivalente es promediar con
una carga de prueba pequeña macroscópicamente pero grande molecularmente, que
promedia sola. \textbf{En el límite de mirar de lejos, el resultado no depende
del detalle del promedio.}
\end{notebox}

\subsection{El campo macroscópico es irrotacional}

No es automático que el campo macroscópico herede las propiedades del
microscópico, así que hay que verificarlo. Tomando el rotacional respecto de
$\vec r$:
\[
\curl \vec E_{\text{macro}}(\vec r) = \int_V dV' \, f(\vec r\,')\,
\curl \vec E_{\text{micro}}(\vec r + \vec r\,') ,
\]
porque $f$ no depende de $\vec r$. Y un desplazamiento del origen no cambia las
derivadas: con $\vec r\,'' = \vec r + \vec r\,'$, el rotacional respecto de
$\vec r$ es el rotacional respecto de $\vec r\,''$, que es cero.

\begin{keybox}[El campo macroscópico deriva de un potencial]
\[
\boxed{\ \curl \vec E_{\text{macro}} = 0
\quad\Longrightarrow\quad
\vec E_{\text{macro}} = -\grad \phi \ }
\]
y también $\oint_C \vec E \cdot d\vec l = 0$ sobre toda curva cerrada.
\end{keybox}

A partir de acá el subíndice se elimina: el único campo que interesa es el
macroscópico. En el vacío ambos coinciden. La ventaja es que el campo
macroscópico es \textbf{suave}, mientras que el microscópico varía sin control en
cuanto uno se mueve a escalas intermedias.

\subsection{La cavidad en forma de aguja}

Falta el paso que permite usar dentro del material la expresión deducida fuera.
Se considera un material \textbf{isótropo} ---sin direcciones privilegiadas--- y
se le practica una \textbf{cavidad en forma de aguja}, pequeña macroscópicamente
pero grande microscópicamente (por ejemplo, una micra de ancho por cien de
largo). La cavidad tiene dos tapas $S_1$, $S_2$ y una superficie lateral $S_L$.

\begin{figure}[H]
  \centering
  % Clase 9 — la cavidad en forma de aguja: el argumento que extiende hacia
% ADENTRO del material un resultado deducido afuera.
% Lo que la figura tiene que mostrar es por que la aguja se elige PARALELA al
% campo: asi la componente tangencial (la unica que la circulacion controla) es
% el campo entero. Y por que sigma_p se anula en la pared lateral: ahi P es
% perpendicular a la normal.
% La aguja va exagerada en grosor para que entren los rotulos; la nota al pie
% aclara la proporcion real, que es 1:100.
\begin{tikzpicture}[scale=1]

  % material dielectrico
  \fill[figblue!8] (-4.2,-1.9) rectangle (4.2,1.9);
  \draw[figline, line width=0.8pt] (-4.2,-1.9) rectangle (4.2,1.9);
  \node[figlbl, figblue] at (-3.4,1.5) {diel\'ectrico};
  \node[figlblaux, align=center] at (-3.35,0.95) {is\'otropo};

  % campo en el dielectrico
  \foreach \y in {-1.35,1.35}{
    \draw[figvecaux] (-2.2,\y) -- (-0.6,\y);}
  \node[figlbl, figgray] at (-1.4,1.72) {$\vec E^{\text{diel}}$};

  % la cavidad (aguja), paralela al campo
  \fill[white] (-1.6,-0.42) rectangle (2.4,0.42);
  \draw[figred, line width=1.0pt] (-1.6,-0.42) rectangle (2.4,0.42);
  \node[figlblaux, figred] at (2.55,1.45) {cavidad (vac\'io)};
  \draw[figaux] (2.25,1.28) -- (1.75,0.5);

  % tapas y lateral
  \node[figlbl, figred] at (-1.95,0) {$S_1$};
  \node[figlbl, figred] at (2.75,0) {$S_2$};
  \node[figlblaux, figred] at (0.4,-0.72) {$S_L$};

  % curva cerrada C
  \draw[figvecres, line width=1.0pt] (-1.2,-0.18) -- (1.9,-0.18);
  \draw[figvecres, line width=1.0pt] (1.9,0.18) -- (-1.2,0.18);
  \node[figlbl, figred, fill=white, inner sep=1.2pt] at (0.35,0.05) {$C$};

  % campo dentro de la cavidad
  \draw[figvec] (0.9,-0.9) -- (2.3,-0.9);
  \node[figlbl, figblue] at (1.6,-1.28) {$\vec E^{\text{cav}}$};

  % P paralela a E, normal lateral perpendicular
  \draw[figvec, figamber] (-3.3,-1.35) -- (-2.2,-1.35);
  \node[figlbl, figamber] at (-2.75,-1.72) {$\vec P$};
  \draw[figvecaux] (0.4,0.42) -- (0.4,0.95);
  \node[figlblaux, figgray] at (0.85,0.72) {$\hat n$};

  \node[figlblaux, align=center] at (0,-2.55)
    {(El grosor est\'a exagerado: la aguja real es del orden de 1 micra de ancho
     por 100 de largo.)};

\end{tikzpicture}

  \caption{Sobre $C$, $\oint \vec E\cdot d\vec l = 0$ da
  $E^{\text{cav}}_t = E^{\text{diel}}_t$; eligiendo la aguja paralela a $\vec E$,
  la tangencial es el campo entero. En $S_L$ se tiene $\hat n \perp \vec P$, así
  que $\sigma_p = 0$: no hay carga concentrada y se puede tomar el límite
  $S_1, S_2, S_L \to 0$.}
\end{figure}

Sobre una curva cerrada $C$ que recorre la aguja a lo largo,
$\oint_C \vec E \cdot d\vec l = 0$, lo que da
$E^{\text{cav}}_t = E^{\text{diel}}_t$. Ahora se elige la cavidad
\textbf{paralela al campo} en el dieléctrico, con lo que la componente tangencial
es el campo entero: $E^{\text{cav}}_t = E^{\text{diel}}$.

\begin{warnbox}[Acá y sólo acá entra la isotropía]
Hasta este punto no se usó. Se usa ahora, y es el paso que cierra el argumento:
como el problema no tiene ninguna otra dirección privilegiada, el campo dentro de
la cavidad no puede apuntar en otra dirección que la del campo del dieléctrico
---si tuviera componente normal, ¿por qué hacia un lado y no hacia el otro?
\end{warnbox}

Además, por isotropía $\vec P$ es paralela al campo, de modo que sobre la
superficie lateral $\sigma_p = \vec P \cdot \hat n = 0$: no hay carga concentrada
en las paredes de la aguja. Sin nada singular, se puede hacer tender $S_1$, $S_2$
y $S_L$ a cero y concluir que $\vec E^{\text{cav}} = \vec E^{\text{diel}}$.

Pero el campo dentro de la cavidad ya se sabe calcular, porque \textbf{está fuera
del material} ---a una distancia macroscópica de la molécula más cercana, aunque
microscópicamente esté pegado.

\begin{keybox}[Conclusión]
Las expresiones de $\phi$ y $\vec E$ deducidas fuera \textbf{valen también
dentro} del dieléctrico, con la restricción de que el material sea isótropo.
\end{keybox}

\section{Ley de Gauss en presencia de dieléctricos}

La ley de Gauss usual sigue valiendo, pero exige \textbf{toda} la carga, libre y
de polarización. Y ahí está el problema: la carga de polarización es en general
una incógnita, se genera sola y no se controla. Se busca entonces una ley de
Gauss escrita sólo en términos de las cargas libres.

\begin{warnbox}[«Carga libre» es un nombre confuso]
No significa que se pueda mover. Una carga depositada sobre un dieléctrico está
fija y es libre. \textbf{Carga libre es la que no es de polarización}; por eso a
veces se prefiere «carga externa», la que uno accede externamente.
\end{warnbox}

Se toma una superficie gaussiana $S$ \textbf{dentro} del dieléctrico, cuyo
material está bordeado por superficies $S_1, S_2, S_3, \dots$. La ley de Gauss da
\[
\oiint_S \vec E \cdot \hat n \, dS = \frac{Q + Q_p}{\varepsilon_0}.
\]

La carga de polarización encerrada se escribe con sus definiciones:
\[
Q_p = \int_V \left(-\div \vec P\right) dV
+ \int_{S_1 \cup S_2 \cup S_3} \vec P \cdot \hat n \, dS ,
\]
donde $V$ es el volumen de material encerrado ---que \textbf{no} incluye lo que
está fuera del dieléctrico--- y las $\hat n$ son salientes del material.
Aplicando el teorema de la divergencia al primer término aparece una integral
sobre \emph{toda} la frontera de $V$, es decir $S \cup S_1 \cup S_2 \cup S_3$, y
las contribuciones de los bordes del material se cancelan contra el segundo
término. Sobrevive únicamente
\[
Q_p = -\oiint_S \vec P \cdot \hat n \, dS .
\]

\begin{notebox}[Eso no es una densidad superficial de verdad]
En la superficie gaussiana no hay nada. Es el sobrante que queda al compensarse
las integrales de los bordes reales del material.
\end{notebox}

Sustituyendo y pasando el término al miembro izquierdo:
\[
\oiint_S \left(\varepsilon_0 \vec E + \vec P\right)\cdot \hat n \, dS = Q .
\]

\begin{keybox}[Vector desplazamiento y ley de Gauss en dieléctricos]
\[
\boxed{\ \vec D = \varepsilon_0 \vec E + \vec P \ }
\]
\[
\boxed{\ \oiint_S \vec D \cdot \hat n \, dS = Q_{\text{libre}}
\qquad\Longleftrightarrow\qquad
\div \vec D = \rho_{\text{libre}} \ }
\]
\end{keybox}

La forma diferencial sale del mismo argumento de volumen pequeño que se usó para
$\vec E$: donde había $\vec E$ ahora hay $\vec D$, y donde había
$Q/\varepsilon_0$ ahora hay $Q$.

\begin{warnbox}[Barrer la mugre bajo la alfombra]
Advertencia explícita del docente: \textbf{no se resolvió nada}. La carga de
polarización molestaba y no se tiró: se escondió dentro de $\vec D$. La
polarización sigue ahí. Es un acto de prestidigitación; estéticamente se avanzó,
pero el problema sigue vigente.
\end{warnbox}

Y el problema es concreto: falta alguien que diga cuánto vale $\vec D$ dado
$\vec E$. Sin eso no se puede volver a $\vec E$, que es el campo que realmente se
observa y el único que deriva de un potencial.

\section{La ecuación constitutiva}

Lo que falta se llama \textbf{ecuación constitutiva} del material: la relación
entre $\vec P$ y $\vec E$ o, equivalentemente, entre $\vec D$ y $\vec E$. Conocer
una es conocer la otra.

Deducirla desde primeros principios es un problema muy difícil. Un pedacito de
materia tiene del orden de $10^{22}$--$10^{23}$ moléculas, y hay que tratar la
interacción de todas ellas con el campo y entre sí. Si cada molécula «habla» sólo
con el campo aplicado y no con las demás, el problema es abordable con
\textbf{mecánica estadística} ---una rama de la física que no se dicta en la
Facultad de Ingeniería---. Si las moléculas interactúan entre sí, es mucho más
difícil.

\begin{warnbox}[Lo que este curso no va a hacer]
El docente dice explícitamente que \textbf{hallar la ecuación constitutiva no se
va a hacer}. Se la toma como dato, medida experimentalmente, y con ella se
resuelven problemas de electrostática ---por ejemplo, un cilindro de material
dentro de un campo eléctrico---. Se usan entonces \textbf{relaciones
constitutivas fenomenológicas}: no deducidas, sino extraídas de experimentos.
\end{warnbox}

\subsection{Susceptibilidad y permitividad}

Se restringe a materiales \textbf{homogéneos}, \textbf{isótropos} y \textbf{sin
polarización espontánea} (si se apaga $\vec E$, se apaga $\vec P$).

\begin{notebox}[Ferroeléctricos y ferromagnéticos]
Los materiales con polarización espontánea existen y se llaman
\textbf{ferroeléctricos}, pero son poco comunes. El fenómeno análogo sí es
frecuentísimo en magnetismo: los imanes conservan la magnetización tras apagar el
campo. Son propiedades distintas: hay muchos ferromagnéticos ---hierro, níquel,
cobalto--- que no son ferroeléctricos.
\end{notebox}

Por isotropía, $\vec P$ no puede apuntar en otra dirección que $\vec E$, y por
homogeneidad la relación es la misma en todo punto. Queda entonces una única
función escalar.

\begin{keybox}[Susceptibilidad y permitividad]
\[
\boxed{\ \vec P = \chi(|\vec E|)\, \vec E \ }
\]
\[
\vec D = \varepsilon_0 \vec E + \vec P
= \left[\varepsilon_0 + \chi(|\vec E|)\right] \vec E
\equiv \varepsilon(|\vec E|)\, \vec E
\]
$\chi$ es la \textbf{susceptibilidad eléctrica} y $\varepsilon$ la
\textbf{permitividad eléctrica} del material.
\end{keybox}

\begin{notebox}[El nombre es adecuado]
Una persona susceptible es la que reacciona mucho ante una perturbación pequeña.
Acá, cuanto mayor es $\chi$, mayor es la polarización para un campo dado.
\end{notebox}

La relación se supone \textbf{local}: la polarización en un punto depende del
campo en ese punto. Existen relaciones no locales, pero son poco comunes y, en
todo caso, el acoplamiento ocurre a distancias microscópicas: a escala
macroscópica la aproximación local es excelente.

\subsection{Por qué tantos materiales son lineales}

Un material es \textbf{lineal} cuando $\chi$ ---o equivalentemente
$\varepsilon$--- no depende de $\vec E$, es decir es una constante:
\[
\vec P = \chi \vec E, \qquad \vec D = \varepsilon \vec E .
\]

Muchísimos materiales realistas lo son, y el motivo se conecta con la advertencia
sobre la caricatura: \textbf{cuando se aplica un campo realista de laboratorio, a
las moléculas casi no les pasa nada}. Siguen desordenadas y oscilando por
agitación térmica, y apenas se alinean una porción minúscula.

Como el efecto es tan pequeño, se puede desarrollar $\chi(|\vec E|)$ en serie de
Taylor alrededor de cero, y la primera corrección resulta de orden $E^2$:
\[
\chi(E) \simeq \chi(0) + \mathcal{O}(E^2),
\]
despreciable porque los campos aplicables en un laboratorio son muy pequeños
comparados con los campos característicos a escala molecular. La susceptibilidad
vale, en la práctica, siempre lo que valdría en el límite $E \to 0$.

\begin{keybox}[Constante dieléctrica]
\[
\boxed{\ K = \frac{\varepsilon}{\varepsilon_0}\ }
\]
\end{keybox}

\begin{warnbox}[Decir «material de constante dieléctrica $K$» ya lo dice todo]
Esa frase sólo tiene sentido si $K$ es efectivamente una constante, y eso exige
que el material sea homogéneo, isótropo, sin polarización espontánea y lineal.
Muchas veces el enunciado de un ejercicio no lista esas hipótesis: al dar $K$,
las está dando implícitamente.
\end{warnbox}

\vspace{1em}
\noindent\emph{La clase siguiente arranca por las consecuencias de esta
definición ---el signo de $\chi$ y la cota $K \ge 1$--- y pasa a los primeros
ejemplos con dieléctricos.}

\end{document}
